\documentclass[letterpaper]{article}

\newif\ifofficialformat
\IfFileExists{aaai2027.sty}{%
  \officialformattrue
  \usepackage[submission]{aaai2027}
}{%
  \officialformatfalse
  \usepackage[margin=1in]{geometry}
  \usepackage{newtxtext}
}

\usepackage{amsmath,amssymb,amsthm,mathtools}
\usepackage{booktabs}
\usepackage{graphicx}
\usepackage{microtype}
\usepackage{url}
\usepackage{xcolor}
\usepackage{enumitem}
\usepackage{multirow}

\graphicspath{{figures/}}

\newtheorem{theorem}{Theorem}
\newtheorem{lemma}[theorem]{Lemma}
\newtheorem{proposition}[theorem]{Proposition}
\newtheorem{corollary}[theorem]{Corollary}
\theoremstyle{definition}
\newtheorem{definition}[theorem]{Definition}
\newtheorem{remark}[theorem]{Remark}

\newcommand{\R}{\mathbb{R}}
\newcommand{\E}{\mathbb{E}}
\newcommand{\Prob}{\mathbb{P}}
\newcommand{\risk}{\mathcal{R}}
\newcommand{\1}{\mathbf{1}}
\newcommand{\asympconst}{\mathrel{\asymp}}
\newcommand{\cO}{\mathcal{O}}
\newcommand{\eps}{\varepsilon}

\title{Beyond Effective Sample Size:\\
Spectral Scaling Laws for Persistent Data Streams}
\author{Anonymous submission}

\begin{document}
\maketitle

\begin{abstract}
Power-law learning curves in linear regression are usually derived from
independent samples, while dependent data are often summarized through a
single effective sample size. We show that a scalar correction misses a
qualitatively different regime: temporal persistence can change the scaling
exponent when it is aligned with the feature spectrum. We introduce
\emph{spectral renewal regression}, a persistent coordinate model with
marginal covariance eigenvalues $\lambda_j\asymp j^{-a}$, mode durations
$\tau_j\asymp j^r$, and target predictive energy
$s_j=\lambda_j\theta_j^2\asymp j^{-b}$. All persistence profiles have the
same one-time covariance. However, statistically new blocks arrive according
to the innovation spectrum $q_j\propto\lambda_j/\tau_j\asymp
j^{-(a+r)}$. For an $M$-coordinate block-averaging estimator trained on $B$
innovation blocks, we derive an exact finite-sample bias--variance identity
and prove
\[
\E\risk(\widehat\theta,\theta)
=\Theta\!\left(
M^{-(b-1)}+B^{-(b-1)/(a+r)}
+\sigma^2\frac{\min\{M,B^{1/(a+r)}\}}{B}
\right).
\]
In the noiseless case, the first two terms are the exact minimax risk over a
signed source class. A renewal-time sandwich further shows that, under
finite-variance dwell times, the same noiseless exponent holds for a trajectory
of exactly $N$ raw observations. Thus uniform persistence ($r=0$) preserves
the usual data exponent, whereas increasingly persistent weak modes ($r>0$)
replace $a$ by $a+r$. We give stationary semi-Markov and reversible Markov
realizations, derive the full model--data scaling function, and verify the
predicted exponents and phase collapse with exact calculations and Monte Carlo
experiments. The result isolates when correlation merely changes
constants and when it changes the rate at which new spectral directions
become learnable.
\end{abstract}

\section{Introduction}

Scaling laws turn small experiments into forecasts of how loss changes with
model size, data, and compute.  Empirical power laws have been observed across
vision, language, speech, and generative modeling
\cite{hestness2017deep,rosenfeld2020constructive,kaplan2020scaling,hoffmann2022training},
while spectral theories explain many such curves through the progressive
resolution of weaker covariance or kernel modes
\cite{sollich1998learning,bordelon2020spectrum,canatar2021spectral,bahri2024explaining,lin2024scaling}.
These results support a useful planning picture: more data reveals increasingly
fine predictive structure, and the unresolved tail determines the learning
slope.

Most scaling-law theories formalize this picture with independent examples.
Yet many high-value datasets are contiguous trajectories: robot experience,
physical and environmental sensors, financial histories, user-interaction
logs, and replay buffers.  A common practical correction replaces the nominal
sample count $N$ by a scalar effective sample size $N_{\rm eff}$ computed from
a global mixing time or integrated autocorrelation.  This is appropriate when
all feature directions share one temporal kernel.  It can be misleading when
the directions that dominate late-stage learning are also the slowest to
refresh.

Consider two streams with exactly the same one-time covariance spectrum.  In
the first, every spectral direction persists for the same duration.  In the
second, high-index directions are individually weak but occur in long bursts.
A global average can assign the two streams the same autocorrelation time.
Nevertheless, late directions arrive much less often in the second stream, so
the number of observations needed to discover them grows with their index.
The relevant object is then not one effective sample count but a sequence of
mode-wise information rates.  This motivates our central question:
\begin{quote}
\emph{When does dependence change a scaling exponent, rather than only a
constant?}
\end{quote}

We answer this question in an exactly solvable persistent regression model.
Mode $j$ occupies a fraction $\lambda_j$ of raw time but lasts $\tau_j$ steps
whenever it appears.  Preserving the same one-time covariance forces a new
block to choose mode $j$ with probability
\begin{equation}
q_j=\frac{\lambda_j/\tau_j}{\sum_\ell\lambda_\ell/\tau_\ell}.
\label{eq:intro-q}
\end{equation}
Thus $\lambda$ is the \emph{marginal spectrum}, while $q$ is the
\emph{innovation spectrum}: it measures the arrival of statistically new
spectral information.  If $\lambda_j\asymp j^{-a}$ and
$\tau_j\asymp j^r$, then $q_j\asymp j^{-(a+r)}$.  After $B$ innovations, the
largest mode seen with constant probability is
$J_B\asymp B^{1/(a+r)}$.  If target energy satisfies
$s_j=\lambda_j\theta_j^2\asymp j^{-b}$, the unresolved tail is
\begin{equation}
\sum_{j>J_B}s_j\asymp B^{-(b-1)/(a+r)}.
\label{eq:intro-mechanism}
\end{equation}
This simple coverage mechanism changes the slope while leaving the marginal
prediction problem unchanged.

Our contributions are as follows.
\begin{enumerate}
\item We define \emph{spectral renewal regression} and derive an exact
finite-sample decomposition into model approximation, unseen-mode bias, and
label-noise variance.  Under power laws, all three terms have matching upper
and lower rates.
\item We prove an exact noiseless minimax equality over a signed source class.
No optimizer, number of epochs, or data-reuse strategy can infer the sign of a
predictive mode that never arrived.
\item We give a direct counterexample to scalar effective sample size.  Two
stationary streams have the same marginal covariance and the same trace
integrated autocorrelation time, yet have exponents $(b-1)/a$ and
$(b-1)/(a+r)$.  We then derive the sample- and compute-optimal consequences of
the mode-wise law.
\item We pass from a fixed number of innovations to exactly $N$ raw examples.
At a renewal boundary, the innovation exponent holds for every $r\ge0$
without a finite-variance dwell assumption.  Under stationary initialization,
a size-biased initial block can create a separate inspection-paradox phase.
\item Using the von Bahr--Esseen inequality, we transfer the noisy hard-target
rate under fractional, rather than second, dwell moments.  This covers
substantial infinite-variance regimes.
\item We prove invariance under known well-conditioned dense representations
and deliberately test the boundary of the theory.  Dense-dictionary
regression retains the predicted slopes, while a dense Gaussian
autoregressive negative control does not: autocorrelation alone is not enough
when every sample contains every mode.
\end{enumerate}

The distinction in the last item is important.  We do not claim that arbitrary
nonseparable time series universally replace $a$ by $a+r$.  The exponent
changes when persistence limits the arrival of new predictive spectral
information.  This narrower claim is both provable and operational: it suggests
estimating the persistence-adjusted spectrum $\lambda_j/\tau_j$, rather than
correcting all directions by one global mixing time.

\section{Related Work}

\paragraph{Scaling laws and spectral learning curves.}
Empirical neural scaling laws were documented by
\cite{kaplan2020scaling} and refined for compute-optimal training by
\cite{hoffmann2022training}. Spectral theories explain learning curves through
resolution of covariance or kernel eigenmodes
\cite{bordelon2020spectrum,bahri2024explaining}. Solvable random-feature
models connect power-law data structure to model- and data-scaling exponents
\cite{maloney2022solvable}, while feature learning can alter those exponents
outside the kernel regime \cite{bordelon2024feature}.

Our closest starting point is the infinite-dimensional sketched linear model
of \cite{lin2024scaling}, which proves a two-resource law for one-pass SGD on
i.i.d. data. \cite{lin2025reuse} adds a source exponent and shows how
multi-pass optimization can improve the iteration term when data are reused.
We instead keep the marginal task fixed and modify the temporal arrival
process. Reusing a persistent block cannot reveal a mode that has not arrived;
this creates an innovation-limited spectral cutoff.

\paragraph{Regression with dependent samples.}
Least squares under Markovian sampling is known to depend on mixing time, and
worst-case minimax rates can be attained by data dropping; structured chains
can benefit from experience replay \cite{nagaraj2020markovian}. Finite-time
SGD analyses and recent single-trajectory lower bounds likewise quantify
Markov dependence through global chain parameters
\cite{doan2020finite,sun2025single}. Recent high-dimensional ridge theory
allows correlated Gaussian examples and derives sharp risk and
cross-validation formulas \cite{atanasov2024correlated}. Its
matrix-Gaussian model has separable covariance
$\E[X_{ti}X_{sj}]=K_{ts}\Sigma_{ij}$: every feature direction inherits the
same sample-correlation matrix. Our construction is nonseparable. Mode $j$
has its own correlation time $\tau_j$, and the alignment between
$\lambda_j$ and $\tau_j$ is precisely what changes the exponent.

\paragraph{Infinite occupancy.}
At innovation times, the problem is an infinite-urn scheme with frequencies
$q_j$. Missing-mass and occupancy asymptotics under regular variation are
well studied \cite{gnedin2007occupancy,benhamou2017concentration}. Our risk is
a \emph{target-weighted} missing mass: unseen mode $j$ costs $s_j$, which may
have a different exponent from $q_j$. We use this weighted structure to derive
the model--data scaling law and minimax regression interpretation.


\section{Spectral Renewal Regression}

Let $(e_j)_{j\ge1}$ be orthonormal. A fresh test example is
\begin{equation}
X=Se_J,\quad \Prob(J=j)=\lambda_j,\quad
\Prob(S=\pm1)=\tfrac12,
\label{eq:test-law}
\end{equation}
where $\lambda_j>0$ and $\sum_j\lambda_j=1$. For target $\theta$, the label is
$Y=\langle X,\theta\rangle+\eps$. Prediction risk is
\begin{equation}
\risk(\widehat\theta,\theta)
=\sum_{j\ge1}\lambda_j(\widehat\theta_j-\theta_j)^2,
\qquad s_j:=\lambda_j\theta_j^2.
\label{eq:risk-source}
\end{equation}
Associate mode $j$ with dwell weight $\tau_j\ge1$, define
\begin{equation}
Z:=\sum_j\lambda_j/\tau_j,\qquad
q_j:=\frac{\lambda_j/\tau_j}{Z},\qquad \mu:=Z^{-1},
\label{eq:q-def}
\end{equation}
and draw $B$ independent innovation blocks $J_k\sim q$. Block $k$ has sign
$S_k$ and provides the average
\begin{align}
\overline Y_k&=S_k\theta_{J_k}+\overline\eps_k,
&\E[\overline\eps_k\mid J_k=j]&=0,\nonumber\\
&&\E[\overline\eps_k^2\mid J_k=j]&=\sigma^2/\tau_j.
\label{eq:block-observation}
\end{align}
For integer $\tau_j$, this is the average of $\tau_j$ repeated noisy labels.
Since $q_j\tau_j/\mu=\lambda_j$, expanding blocks into raw observations leaves
the one-time covariance unchanged.

Let $K_j=\sum_{k\le B}\1\{J_k=j\}$. For a model restricted to the first $M$
modes, use
\begin{equation}
\widehat\theta_j=
\begin{cases}
K_j^{-1}\sum_{k:J_k=j}S_k\overline Y_k,&K_j>0,\\
0,&K_j=0,
\end{cases}
\quad j\le M,
\label{eq:estimator}
\end{equation}
with $\widehat\theta_j=0$ for $j>M$.

\paragraph{Dense representations.}
The coordinate notation is a latent basis, not a requirement that observed
vectors be sparse. Let $V\in\R^{D\times D}$ be invertible with rows $v_j^\top$,
and observe $X=Sv_J$. Labels are $S\theta_J$, so the ambient teacher is
$w^\star=V^{-1}\theta$.

\begin{proposition}[Representation invariance]
\label{prop:dense}
For every latent estimate $\widehat\theta$, the ambient estimate
$\widehat w=V^{-1}\widehat\theta$ has exactly the same risk:
\begin{equation}
\E\langle X,\widehat w-w^\star\rangle^2
=\sum_{j=1}^D\lambda_j(\widehat\theta_j-\theta_j)^2.
\label{eq:dense-risk}
\end{equation}
If $\lambda_1\ge\cdots\ge\lambda_D$ and
$mI\preceq VV^\top\preceq LI$, the ordered eigenvalues of
$\Sigma=V^\top\operatorname{diag}(\lambda)V$ obey
$m\lambda_j\le\lambda_j(\Sigma)\le L\lambda_j$.
\end{proposition}
The proof is a change of coordinates plus the min--max principle. Thus a known,
uniformly well-conditioned dense representation preserves both exact risk and
the marginal spectral exponent.

\section{Exact Risk and Scaling Law}

For $K\sim\mathrm{Binomial}(B,q)$ define
\begin{equation}
h_B(q):=\E\left[\frac{\1\{K>0\}}{K}\right].
\label{eq:h-def}
\end{equation}

\begin{theorem}[Exact finite-block risk]
\label{thm:exact-risk}
For every $B,M\ge1$ and every target with
$\sum_j s_j<\infty$, the estimator in Equation~\eqref{eq:estimator} satisfies
\begin{align}
\E\risk(\widehat\theta,\theta)
={}&\underbrace{\sum_{j>M}s_j}_{\text{model approximation}}
+\underbrace{\sum_{j\le M}s_j(1-q_j)^B}_{\text{unseen-mode bias}}
\nonumber\\
&+\underbrace{\sigma^2\sum_{j\le M}
\frac{\lambda_j}{\tau_j}h_B(q_j)}_{\text{label-noise variance}}.
\label{eq:exact-risk}
\end{align}
\end{theorem}

\begin{proof}
If $K_j=0$, coordinate $j\le M$ is estimated by zero and contributes
$s_j$. If $K_j=k>0$, the signed average is unbiased and has conditional
variance $\sigma^2/(\tau_jk)$. Since $K_j\sim\mathrm{Binomial}(B,q_j)$,
averaging these cases and adding the unavoidable $j>M$ tail gives
Equation~\eqref{eq:exact-risk}.
\end{proof}

The variance term is governed by a simple reciprocal-count estimate.

\begin{lemma}[Reciprocal binomial count]
\label{lem:binomial}
There are universal constants $0<c<C<\infty$ such that, for all $B\ge1$ and
$q\in(0,1]$,
\begin{equation}
c\min\{Bq^2,B^{-1}\}
\le qh_B(q)
\le C\min\{Bq^2,B^{-1}\}.
\label{eq:binomial-bound}
\end{equation}
\end{lemma}

The proof is given in the technical supplement. We now state the principal scaling
law. The notation $u_j\asympconst v_j$ means that their ratio is bounded above
and below by positive constants independent of $j$, $B$, and $M$.

\begin{theorem}[Persistent spectral scaling law]
\label{thm:scaling}
Assume
\begin{equation}
\lambda_j\asympconst j^{-a},
\qquad
\tau_j\asympconst j^r,
\qquad
s_j\asympconst j^{-b},
\label{eq:power-laws}
\end{equation}
where $a>1$, $r\ge0$, and $b>1$. Put $p=a+r$. Then, uniformly for
$B,M\ge2$,
\begin{equation}
\boxed{
\E\risk(\widehat\theta,\theta)
\asympconst
M^{-(b-1)}
+B^{-(b-1)/p}
+\sigma^2\frac{\min\{M,B^{1/p}\}}{B}.}
\label{eq:main-scaling}
\end{equation}
The hidden constants depend only on the comparison constants in
Equation~\eqref{eq:power-laws} and on $a,b,r$.
\end{theorem}

\paragraph{Interpretation.}
The innovation probabilities satisfy $q_j\asymp j^{-p}$. A mode has a
constant probability of appearing once $Bq_j\gtrsim1$, giving the spectral
resolution
\begin{equation}
J_B\asymp B^{1/p}=B^{1/(a+r)}.
\label{eq:JB}
\end{equation}
The first two terms in Equation~\eqref{eq:main-scaling} are equivalently
$\Theta(\min\{M,J_B\}^{-(b-1)})$. The third term counts the number of
well-observed modes, up to the same resolution, divided by $B$.

\begin{corollary}[When persistence changes the exponent]
\label{cor:exponent}
In the noiseless case, or for fixed $\sigma^2=O(1)$ in the hard-target regime
$1<b<a+r$ where the bias term dominates the variance term,
\begin{equation}
\E\risk(\widehat\theta,\theta)
\asympconst M^{-(b-1)}+B^{-(b-1)/(a+r)}.
\label{eq:clean-law}
\end{equation}
Uniform persistence has $r=0$ and preserves the i.i.d. exponent
$(b-1)/a$. If weaker modes persist longer with $r>0$, the data exponent becomes
$(b-1)/(a+r)$.
\end{corollary}

\begin{corollary}[Noise regimes and model choice]
\label{cor:model-choice}
Let $\sigma>0$ be fixed. If $b<a+r$, taking
$M\asymp B^{1/(a+r)}$ achieves risk
$\Theta(B^{-(b-1)/(a+r)})$. If $b>a+r$, the rate-optimal model size is
\begin{equation}
M_\star\asymp(B/\sigma^2)^{1/b},
\label{eq:optimal-M}
\end{equation}
and the resulting dominant risk is
$\Theta((\sigma^2/B)^{(b-1)/b})$.
\end{corollary}


\section{An Exact Information-Theoretic Lower Bound}

The changed exponent is not caused by an inefficient learning algorithm. It
is the exact minimax rate for a natural source class.

Fix positive energies $(s_j)$ and define the signed class
\begin{equation}
\Theta_s
:=\left\{\theta^\omega:
\theta_j^\omega=\omega_j\sqrt{s_j/\lambda_j},\;
\omega_j\in\{-1,+1\}\right\}.
\label{eq:signed-class}
\end{equation}
Let $\mathcal A_M$ be the estimators whose output is supported on
$\{1,\ldots,M\}$.

\begin{theorem}[Exact noiseless minimax risk]
\label{thm:minimax}
When $\sigma=0$,
\begin{equation}
\inf_{\widetilde\theta\in\mathcal A_M}
\sup_{\theta\in\Theta_s}
\E\risk(\widetilde\theta,\theta)
=
\sum_{j>M}s_j+\sum_{j\le M}s_j(1-q_j)^B.
\label{eq:minimax}
\end{equation}
The estimator in Equation~\eqref{eq:estimator} attains the equality.
\end{theorem}

\begin{proof}
For the upper bound, an observed noiseless block reveals
$S_k\overline Y_k=\theta_{J_k}$ exactly, so
Equation~\eqref{eq:estimator} incurs $s_j$ precisely when $j>M$ or when
$j\le M$ is unseen. Its risk is independent of the sign vector $\omega$.

For the lower bound, put the product Rademacher prior on $\omega$. If mode
$j\le M$ is unseen, the data are independent of $\omega_j$, the posterior
remains symmetric, and the Bayes squared prediction error in that coordinate
is at least $s_j$. This event has probability $(1-q_j)^B$. Coordinates
$j>M$ incur $s_j$ because of the model restriction. Summing gives the
right-hand side of Equation~\eqref{eq:minimax}; Bayes risk lower-bounds
minimax risk.
\end{proof}

The theorem clarifies the role of data reuse. Repeated computation on the same
$B$ blocks can reduce optimization error but cannot identify the sign of an
unseen mode. The innovation spectrum sets an irreducible data-coverage floor.


\section{Stationary Correlated-Stream Realizations}

The block experiment has a direct stationary interpretation. Suppose first
that integer durations $\ell_j\asymp j^r$ are used. Concatenate i.i.d. blocks
with modes drawn from $q$ and durations $\ell_j$. Initialize the block
containing time zero from the equilibrium size-biased distribution and choose
its age uniformly within the block. This produces a stationary semi-Markov
stream. Equation~\eqref{eq:occupancy} shows that its one-time covariance is
$\sum_j\lambda_je_je_j^\top$.

Let $N_B=\sum_{k=1}^B\ell_{J_k}$ be the raw duration of $B$ completed blocks.
Then
\begin{equation}
\E N_B=\mu B,
\qquad
N_B/B\longrightarrow\mu\quad\text{almost surely}.
\label{eq:renewal-slln}
\end{equation}
Consequently, indexing by expected raw observations changes constants but not
any exponent. If $r<a-1$, the duration has finite variance and $N_B$ has
ordinary square-root concentration. We keep Theorem~\ref{thm:exact-risk}
indexed by completed innovations because it yields an exact finite-sample
identity and avoids boundary effects from a partially observed final block.
Almost-sure equivalence by itself does not imply an expected-risk equivalence:
rare, unusually long blocks could dominate. The following result controls that
possibility directly.

\begin{theorem}[Fixed raw horizon and the inspection paradox]
\label{thm:raw-horizon}
Assume Equation~\eqref{eq:power-laws}, $\sigma=0$, integer dwell times
$\tau_j\ge1$, and
\begin{equation}
1<b<a+r+1.
\label{eq:raw-horizon-conditions}
\end{equation}
First start the concatenated stream at a renewal boundary and observe exactly
$N$ raw examples.  For every $r\ge0$, uniformly for $N,M\ge2$,
\begin{equation}
\E\risk(\widehat\theta^{(N)},\theta)
\asympconst
M^{-(b-1)}+N^{-(b-1)/(a+r)}.
\label{eq:raw-horizon-rate}
\end{equation}
No finite-variance assumption on the dwell time is needed.

Now initialize the stream from its stationary equilibrium renewal law.  If
$r>0$, then
\begin{equation}
\E\risk(\widehat\theta^{(N)}_{\rm eq},\theta)
\asympconst
M^{-(b-1)}+N^{-(b-1)/(a+r)}+N^{-(a-1)/r}.
\label{eq:stationary-raw-rate}
\end{equation}
For $r=0$, the last term is absent.  Thus equilibrium initialization preserves
the boundary-start exponent when $(b-1)/(a+r)<(a-1)/r$, but otherwise the
forward recurrence time creates a slower inspection-paradox phase.
\end{theorem}

The proof follows individual mode-hitting times rather than concentrating the
total number of renewals.  If $H_j$ is the raw start time of the first block in
mode $j$, then
\begin{equation}
\E H_j=\mu/q_j-\tau_j\le\mu/q_j,
\qquad
(1-q_j)^N\le\Prob(H_j\ge N)
\le\min\{1,\mu/(Nq_j)\}.
\label{eq:hitting-summary}
\end{equation}
Summing these inequalities against $s_j$ gives
Equation~\eqref{eq:raw-horizon-rate} for all $r$.  Under stationary
initialization, the residual duration $R$ of the initial size-biased block has
\begin{equation}
\Prob(R\ge N)\asympconst N^{-(a-1)/r},
\label{eq:inspection-tail-main}
\end{equation}
which yields the additional term in Equation~\eqref{eq:stationary-raw-rate}.
The technical supplement gives matching upper and lower bounds.

A Markov realization makes the nonseparable covariance explicit. Consider
states $(j,s)$ with $s\in\{-1,+1\}$ and transition kernel
\begin{align}
P((j,s),(k,t))
={}&\left(1-\tau_j^{-1}\right)
\1\{(j,s)=(k,t)\}
+\tau_j^{-1}\frac{q_k}{2}.
\label{eq:markov-kernel}
\end{align}

\begin{proposition}[Reversible mode-dependent persistence]
\label{prop:markov}
The chain in Equation~\eqref{eq:markov-kernel} is reversible with stationary
law $\pi_{j,s}=\lambda_j/2$. For $X_t=S_te_{J_t}$,
\begin{equation}
\E[X_tX_{t+\ell}^\top]
=\sum_{j\ge1}\lambda_j
\left(1-\tau_j^{-1}\right)^\ell e_je_j^\top.
\label{eq:mode-autocovariance}
\end{equation}
Unless all $\tau_j$ are equal, this covariance cannot be written as a scalar
time kernel multiplied by the spatial covariance.
\end{proposition}

\begin{proof}
For distinct states,
\[
\pi_{j,s}P((j,s),(k,t))
=\frac{\lambda_jq_k}{4\tau_j}
=\frac{Zq_jq_k}{4},
\]
which is symmetric in $(j,s)$ and $(k,t)$; self-transitions are immediate.
For the covariance, the only contribution correlated with $X_t$ is the event
that no refresh occurs during the next $\ell$ steps. A refresh resamples a
mean-zero sign, giving Equation~\eqref{eq:mode-autocovariance}.
\end{proof}

The proposition explains why a single sample-correlation matrix is
insufficient here: each spatial mode carries a different temporal kernel.
The exact noisy theorem uses deterministic or block-averaged durations;
extending its sharp variance constants to geometric dwell times is discussed
under limitations.


\section{Exact and Monte Carlo Experiments}

All experiments use direct summation or sampling and run on a laptop. Unless
stated otherwise, we set $a=2$, $b=1.8$, normalize each power law by its zeta
constant, and use $M=2^{18}$. The code evaluates
\begin{equation}
\risk_{B,M}^{(0)}
=\sum_{j>M}s_j+\sum_{j\le M}s_j(1-q_j)^B
\label{eq:exact-experiment}
\end{equation}
without fitting any free constants.

\paragraph{Exponent recovery.}
Figure~\ref{fig:learning-curves} varies the persistence exponent. The predicted
slopes are $-0.4$, $-1/3$, and $-2/7$ for $r=0,0.4,0.8$, respectively. Fits
over the seven largest data scales give $-0.399999$, $-0.333333$, and
$-0.285714$. The marginal covariance spectrum is identical across all three
curves; only the innovation spectrum changes.

\begin{figure}[t]
\centering
\includegraphics[width=\columnwidth]{learning_curves.pdf}
\caption{Exact finite-$B$ learning curves. Increasing mode-dependent
persistence changes the asymptotic slope from
$-(b-1)/a$ to $-(b-1)/(a+r)$.}
\label{fig:learning-curves}
\end{figure}

\paragraph{Model--data phase collapse.}
Theorem~\ref{thm:scaling} predicts that the noiseless risk depends on the
effective resolution $\min\{M,B^{1/(a+r)}\}$. Figure~\ref{fig:collapse}
varies both $M$ and $r$. Curves approach the common slope $-(b-1)$ when
plotted against this resolution, supporting the predicted crossover frontier
$M\asymp B^{1/(a+r)}$.

\begin{figure}[t]
\centering
\includegraphics[width=\columnwidth]{phase_collapse.pdf}
\caption{Model--data phase collapse under the effective spectral resolution
$\min\{M,B^{1/(a+r)}\}$. The dashed line has slope $-(b-1)$.}
\label{fig:collapse}
\end{figure}

\paragraph{Sampling validation.}
We independently sample $J_k$ from the infinite Zipf innovation law, construct
the coordinate-averaging estimator, and average its conditional risk over
400 datasets. Across the tested scales and persistence values, the largest
relative deviation from Equation~\eqref{eq:exact-experiment} is $2.1\%$; all
curves agree within Monte Carlo uncertainty. A second experiment samples
integer-duration streams stopped after exactly $N$ raw observations. Its fitted
slopes $-0.399$, $-0.339$, and $-0.290$ closely match the predictions
$-0.400$, $-0.333$, and $-0.286$. The technical supplement provides both plots
and implementation details.


\section{Discussion and Limitations}

\paragraph{What the result establishes.}
Two streams can have the same marginal covariance spectrum and the same
expected number of raw observations, yet different learning exponents. The
relevant object is the innovation spectrum $\lambda_j/\tau_j$. A scalar
effective sample size remains adequate under uniform persistence, but
mode-dependent persistence creates a hierarchy of effective sample sizes and
moves the model--data frontier.

\paragraph{What remains open.}
First, the fixed-raw-horizon theorem is noiseless and assumes finite-variance
dwell times. Noisy stopping rules and the heavy-dwell regime $r\ge a-1$ can
exhibit additional renewal fluctuations and remain open. Second, the
coordinate model isolates coverage and is intentionally sparse. Dense Gaussian
or random feature designs may combine this mechanism with sample-covariance
concentration; proving universality beyond spectral sampling is an important
next step. Third, the clean noisy identity uses deterministic replicated
blocks or block averages. Geometric dwell lengths can add finite-block
corrections to inverse run-length moments. Finally, current experiments are
synthetic. A full empirical study should estimate mode-wise persistence on
real sequential regression datasets and test whether the innovation spectrum
predicts out-of-range learning curves better than a global correlation time.

\paragraph{Practical implication.}
The result suggests measuring not only how much variance lies in a feature
direction, but how often that direction is refreshed. When weak directions
arrive in long bursts, collecting a longer contiguous trajectory can be less
valuable than increasing temporal or environmental diversity, even when the
number of logged observations is unchanged.


\section{Conclusion}

We introduced an exactly solvable regression model in which temporal
persistence is aligned with a power-law feature spectrum. The one-time
covariance remains fixed, but new spectral information arrives according to
$q_j\propto\lambda_j/\tau_j$. This changes the learnable resolution from
$B^{1/a}$ to $B^{1/(a+r)}$ and the hard-target data exponent from
$(b-1)/a$ to $(b-1)/(a+r)$. Exact finite-sample risk, a matching noiseless
minimax theorem, a fixed-raw-horizon transfer, and stationary process
realizations make clear that correlation can do more than rescale sample size:
it can alter which spectral modes the data ever reveal.


\ifofficialformat
  \bibliographystyle{aaai2027}
\else
  \bibliographystyle{plain}
\fi
\bibliography{references}

\end{document}
